% SHARD-ID: AQM-60-PREQUANTISATION-POLARISATION-DERHAM-CSS
% SHARD-TITLE: Prequantisation, Polarisation, and the de Rham CSS Test
% SHARD-SUMMARY: Separates closed-point Weyl prequantisation, tangent-cotangent polarisation, and the finite de Rham CSS test.
% SHARD-SUMMARY: Proves that tangent and cotangent spaces give natural local Lagrangians but not a dynamical principle by themselves.
% SHARD-SUMMARY: Works out product-ring, finite-field circle, and formal real-line examples illustrating the distinction.
% SHARD-KEYWORDS: prequantisation, polarisation, tangent, cotangent, Weyl, de Rham CSS, circle, real line

\section{Prequantisation, Polarisation, and the de Rham CSS Test}
\label{sec:prequantisation-polarisation-derham-css}

\subsection{Status and Source Anchors}

\begin{sloppypar}
This shard corrects the interpretation of AQM-54--AQM-59.  The residue-field
Weyl layer is AQM-28 and convention (u).  The tangent and cotangent conventions
are convention (ao), with K\"ahler differentials sourced from
\path{references/algebraic_geometry/stacks_project_algebra.tex}, lines
33793--33872 and 34310--34340, and tangent spaces sourced from
\path{references/algebraic_geometry/stacks_project_varieties.tex}, lines
2915--3019.  The stabilizer/CSS isotropy dictionary is AQM-08, sourced there
from Gottesman, Ashikhmin--Knill, and Bombin--Martin--Delgado.  The real Weyl
comparison is AQM-19.  The finite-field circle calculation is AQM-59.
\end{sloppypar}

The conclusion is deliberately conservative: tangent-cotangent geometry
provides natural local polarisations of a geometric phase space.  It does not
by itself provide dynamics.  The finite \(\texttt{de\_rham\_css\_test}\) is a
separate chain-complex rule, useful as a calculable test but not identical to
the tangent-cotangent polarisation picture.

\subsection{Lamport Dependency Skeleton}

\[
\begin{array}{rcl}
D1&:&k\text{ is a field, }X=\Spec A,\ P\in X\text{ has residue field }
      K=\kappa(P).\\
D2&:&E_P^{\rm res}=K_q\oplus K_p,\quad
      \omega_{\rm res}((q,p),(q',p'))=pq'-qp'.\\
D3&:&\mathcal H_P^{\rm res}=\ell^2(K),\text{ with the finite Weyl operators
      of AQM-28 when }K\text{ is finite}.\\
D4&:&T_P^*X=\Omega^1_{A/k}\otimes_A K,\quad
      T_PX=\operatorname{Hom}_K(T_P^*X,K).\\
D5&:&E_P^{\rm geom}=T_PX\oplus T_P^*X.\\
D6&:&\omega_{\rm geom}((v,\alpha),(w,\beta))=\beta(v)-\alpha(w).\\
L1&:&E_P^{\rm geom}\text{ is symplectic when }T_PX\text{ is finite
      dimensional}.\\
L2&:&T_PX\oplus0\text{ and }0\oplus T_P^*X\text{ are Lagrangian subspaces}.\\
D7&:&\text{For finite }K,\text{ a geometric Weyl representation may be placed
      on }\ell^2(T_PX).\\
D8&:&C^0=A\to C^1=\Omega^1_{A/k}\to C^2=\Omega^2_{A/k}.\\
D9&:&R_X=\operatorname{im}d_0,\quad R_Z=\operatorname{im}d_1^T.\\
T1&:&D2\text{--}D3,\ D4\text{--}D7,\text{ and }D8\text{ are distinct layers}.\\
T2&:&\text{A polarisation is not yet dynamics; extra global data are required.}
\end{array}
\]

\subsection{Local Phase Space and Polarisation D4--L2}

Let \(V=T_PX\) and \(V^*=T_P^*X\).  Define
\[
  E_P^{\rm geom}=V\oplus V^*
\]
and
\[
  \omega_{\rm geom}((v,\alpha),(w,\beta))=\beta(v)-\alpha(w).
\]
The form is bilinear and alternating because
\[
  \omega_{\rm geom}((v,\alpha),(v,\alpha))=\alpha(v)-\alpha(v)=0.
\]
It is nondegenerate: if \((v,\alpha)\ne0\), then either \(v\ne0\) or
\(\alpha\ne0\).  If \(v\ne0\), choose \(\beta\in V^*\) with \(\beta(v)\ne0\);
then \(\omega_{\rm geom}((v,\alpha),(0,\beta))=\beta(v)\ne0\).  If
\(\alpha\ne0\), choose \(w\in V\) with \(\alpha(w)\ne0\); then
\(\omega_{\rm geom}((v,\alpha),(w,0))=-\alpha(w)\ne0\).  This proves \(L1\).

The subspace \(V\oplus0\) is isotropic because both covector entries are zero:
\[
  \omega_{\rm geom}((v,0),(w,0))=0.
\]
It has dimension \(\dim_KV=\frac12\dim_KE_P^{\rm geom}\), so it is Lagrangian.
The same proof gives that \(0\oplus V^*\) is Lagrangian.  These are the two
standard local polarisations.

For finite \(K\), the usual finite Weyl model for \(E_P^{\rm geom}\) is
\[
  \mathcal H_P^{\rm geom}=\ell^2(V),
\]
with
\[
  T(v_0)|v\rangle=|v+v_0\rangle,\qquad
  R(\alpha)|v\rangle=\psi_K(\alpha(v))|v\rangle,
\]
where \(\psi_K\) is the additive character used in AQM-28.  If
\(\dim_KV=m\), then
\[
  \ell^2(V)\cong \ell^2(K^m)\cong \ell^2(K)^{\otimes m}
\]
after choosing a basis of \(V\).  Thus a singular or nonreduced point with
\(m>1\) gives \(m\) local \(K\)-qudits in the geometric tangent-cotangent model,
not one.

\subsection{Why This Is Not the Same as the de Rham CSS Test}

The residue-Weyl prequantisation is
\[
  E_P^{\rm res}=K_q\oplus K_p,\qquad \mathcal H_P^{\rm res}=\ell^2(K).
\]
This uses only the residue field.  It does not see whether \(P\) is smooth,
singular, reduced, nonreduced, tangent-rich, or tangent-zero.  It gives a local
Heisenberg system but no distinguished isotropic subspace.

The tangent-cotangent construction instead uses \(T_PX\) and \(T_P^*X\).  It
agrees with \(K^2\) only after extra choices when \(\dim_KT_PX=1\).  It can
vanish, as for a reduced zero-dimensional point, or grow, as for singular and
nonreduced points.

The finite \(\texttt{de\_rham\_css\_test}\) is again different.  It uses the
global algebraic complex
\[
  C^0=A\xrightarrow{d_0}C^1=\Omega^1_{A/k}
       \xrightarrow{d_1}C^2=\Omega^2_{A/k}
\]
when these are finite-dimensional \(k\)-vector spaces, puts physical qudits on
a chosen basis of \(C^1\), and forms
\[
  R_X=\operatorname{im}d_0,\qquad R_Z=\operatorname{im}d_1^T.
\]
The equation \(d_1d_0=0\) proves CSS isotropy.  This rule is a chain-complex
stabilizer factory.  It is not the same statement as ``at each point,
\(T_PX\oplus0\) and \(0\oplus T_P^*X\) are Lagrangian.''

\subsection{Example I: Product Fields}

Let
\[
  A=k^n=k e_1\oplus\cdots\oplus k e_n .
\]
AQM-58 proves that every \(k\)-derivation \(D:A\to M\) vanishes.  Hence
\[
  \Omega^1_{A/k}=0.
\]
At the \(i\)-th point \(P_i\), the tangent and cotangent spaces are
\[
  T_{P_i}^*X=0,\qquad T_{P_i}X=0.
\]
Thus
\[
  E_{P_i}^{\rm geom}=0.
\]
There is no geometric infinitesimal mode to polarise.

The residue-Weyl layer is different:
\[
  E_{P_i}^{\rm res}=k^2,\qquad \mathcal H_{P_i}^{\rm res}=\ell^2(k).
\]
For the full product ring this gives the \(n\)-qudit ambient Hilbert space
\[
  \mathcal H^{\rm res}=\ell^2(k)^{\otimes n}.
\]
Therefore the product-field example is the clean warning case.  Residue
prequantisation supplies kinematics, but tangent-cotangent geometry supplies no
nonzero local phase space, and the de Rham CSS rule supplies no stabilizer.

\subsection{Example II: The Odd Finite-Field Circle}

Let \(k=\mathbb F_q\) with \(q\) odd and
\[
  C=\Spec k[x,y]/(x^2+y^2-1).
\]
For a rational point \(P=(a,b)\), AQM-59 computes
\[
  T_P^*C=(k\,dx\oplus k\,dy)/(a\,dx+b\,dy).
\]
The tangent space is the kernel of the differential of the equation:
\[
  T_PC=\{(u,v)\in k^2:au+bv=0\}.
\]
Since \(a^2+b^2=1\), the pair \((a,b)\) is nonzero, so both spaces are
one-dimensional.

If \(b\ne0\), choose
\[
  \tau_P=(1,-a/b)\in T_PC,\qquad \eta_P=dx\in T_P^*C .
\]
Then \(\eta_P(\tau_P)=1\).  If \(b=0\), then \(a=\pm1\), and choose
\[
  \tau_P=(0,1),\qquad \eta_P=dy,
\]
again with \(\eta_P(\tau_P)=1\).  Therefore
\[
  E_P^{\rm geom}=k\tau_P\oplus k\eta_P
\]
with
\[
  \omega_{\rm geom}((r\tau_P,s\eta_P),(r'\tau_P,s'\eta_P))
  =s'r-sr'.
\]
After this basis choice, \(E_P^{\rm geom}\cong k^2\), so for a smooth rational
point it matches the residue-Weyl label shape.  The match depends on choosing
the tangent basis.

For \(q=3\),
\[
  C(\mathbb F_3)=\{(1,0),(2,0),(0,1),(0,2)\}.
\]
The local basis choices are
\[
\begin{array}{c|c|c}
P & \tau_P & \eta_P\\ \hline
(1,0)&(0,1)&dy\\
(2,0)&(0,1)&dy\\
(0,1)&(1,0)&dx\\
(0,2)&(1,0)&dx
\end{array}
\]
Using the position polarisation \(T_PC\oplus0\) as stabilizer support gives
the local \(X\)-type labels
\[
  \langle X_{(1,0)},X_{(2,0)},X_{(0,1)},X_{(0,2)}\rangle .
\]
Using the momentum polarisation \(0\oplus T_P^*C\) gives instead
\[
  \langle Z_{(1,0)},Z_{(2,0)},Z_{(0,1)},Z_{(0,2)}\rangle .
\]
Either full polarisation produces a product polarisation state after \(+1\)
phases are chosen.  It does not couple different points and does not define a
nontrivial code.  This is not a failure of symplectic geometry; it says that a
polarisation alone is not a dynamical principle.

\subsection{Example III: The Formal Real Affine Line}

Let
\[
  A=\mathbb R[x],\qquad X=\Spec A.
\]
AQM-58 computes
\[
  \Omega^1_{A/\mathbb R}=A\,dx.
\]
At the real point \(P_a=(x-a)\),
\[
  T_{P_a}^*X=\mathbb R\,dx_a,\qquad T_{P_a}X=\mathbb R\,\partial_x|_a .
\]
The geometric local phase space is
\[
  E_{P_a}^{\rm geom}=
  \mathbb R\,\partial_x|_a\oplus \mathbb R\,dx_a,
\]
with
\[
  \omega((r\partial_x,s\,dx),(r'\partial_x,s'\,dx))=s'r-sr'.
\]
The two local polarisations are the tangent line and the cotangent line.  The
global formal cotangent bundle is
\[
  T^*X=\Spec\mathbb R[x,p],
\]
and the usual formal Schr\"odinger quantisation uses \(x\) by multiplication
and \(p=-i\,d/dx\) on \(L^2(\mathbb R)\).  The cotangent calculation explains
the phase-space variables and polarisations.  It does not choose the free
particle Hamiltonian \(p^2/2m\), a potential \(V(x)\), a domain, or an equation
of motion.  Those are further dynamical choices.

\subsection{Operational Conclusion}

The corrected workflow is:
\[
\begin{array}{ll}
1. & \text{Build the closed-point Weyl prequantisation from residue fields.}\\
2. & \text{Compute }T_PX\text{ and }T_P^*X\text{ to get geometric phase spaces.}\\
3. & \text{Record natural local polarisations, but do not call them dynamics.}\\
4. & \text{Add a separate global principle if one wants stabilizer codes,
     Hamiltonians, or propagation.}
\end{array}
\]
The finite \(\texttt{de\_rham\_css\_test}\) remains useful, but it should be
read as one proposed finite chain-complex stabilizer mechanism, not as the
canonical meaning of tangent-cotangent polarisation.
